% Source: OpenAI, "A two-dimensional area law from a global spectral gap"
% (September 24, 2026, openai/math@adc7f1241b42e322a6451854ab7e4b4c146bf78a),
% build/sections/03-quasilocal.tex, lines 14--41 and 52--128 (Lemma 4.1,
% label lem:quasilocal-lr). The general graph estimates are in the chapter on
% local quantum circuits; this section specializes them to lattice domains.

\section{Graph-distance propagation in lattice domains}

Let $\Lambda\subseteq\Z^2$ be a finite domain with induced graph distance
$d_\Lambda$, and let $H=\sum_Xh_X$ be a finite-range Hamiltonian on $\Lambda$
with range $R$ and term-norm bound $J$, as in
Definition~\ref{def:area_law_hamiltonian}. Put $v_R=1+2R(R+1)$ and
$\mu_R=2^{v_R-1}$, let $\tau_t(B)=e^{itH}Be^{-itH}$, and let
$N_l(a)=\{x\in\Lambda:d_\Lambda(a,x)\le l\}$. For a set $K$ of sites, $E_K$ is
the normalized partial trace of Definition~\ref{def:qc_site_expectation}.

\begin{lemma}[Lattice-domain counts]
    \label{lem:area_domain_graph_counts}
    \lean{TNLean.PEPS.AreaLaw.latticeL1Distance_le_one_of_domainGraph_adj,
        TNLean.PEPS.AreaLaw.card_domainGraph_sphere_le,
        TNLean.PEPS.AreaLaw.IsAdmissibleSupport.edist_le,
        TNLean.PEPS.AreaLaw.IsAdmissibleSupport.card_le,
        TNLean.PEPS.AreaLaw.LocalHamiltonian.nonneg_of_support,
        TNLean.PEPS.AreaLaw.LocalHamiltonian.sum_norm_term_containing_le}
    \leanok
    \uses{def:area_law_hamiltonian,thm:area_diamond_count,
        thm:area_diamond_budget}
    Adjacent sites of $\Lambda$ have ambient displacement one, and at most
    $2(d+1)^2$ sites lie at distance exactly $d$ from any site. Every
    original support has $d_\Lambda$-diameter at most $R$ and at most $v_R$
    sites, and for every site $x$,
    \begin{align}
        \sum_{X\ni x}\|h_X\| & \le\mu_RJ.
        \label{eq:area_domain_site_budget}
    \end{align}
\end{lemma}
\begin{proof}
    \leanok
    \uses{thm:area_diamond_count,thm:area_diamond_budget}
    A nearest-neighbor edge changes one coordinate by one. The sphere of
    radius $d$ lies in the diamond of radius $d$, which has
    $1+2d(d+1)\le2(d+1)^2$ sites. The remaining bounds are the diamond
    support and multiplicity counts applied to the family of original
    supports, with every term norm at most $J$; one support exists, so
    $J\ge0$.
\end{proof}

\begin{lemma}[Graph-distance propagation]
    \label{lem:quasilocal-lr}
    \lean{TNLean.PEPS.AreaLaw.exists_quasilocal_lieb_robinson,
        TNLean.PEPS.AreaLaw.LocalHamiltonian.norm_commutator_le_graphDistance,
        TNLean.PEPS.AreaLaw.LocalHamiltonian.commutator_eq_zero_of_edist_eq_top,
        TNLean.PEPS.AreaLaw.LocalHamiltonian.norm_sub_siteExpectation_graphBall_le}
    \leanok
    \uses{def:area_law_hamiltonian,lem:area_domain_graph_counts,
        thm:qc_graph_lieb_robinson,thm:qc_graph_lieb_robinson_zero,
        thm:qc_graph_ball_localization,def:qc_site_expectation}
    Let $q\ge1$. There are positive constants $C,v,c$, depending only on
    $q,R,J$, with the following property. For every domain $\Lambda$, every
    Hamiltonian $H$ as above, every original support $X$ with anchor
    $a\in X$, every operator $A$ supported in $X$, every site $y$, every
    on-site operator $B_y$ at $y$, and every real $t$,
    \begin{align}
        \|[\tau_t(A),B_y]\| & \le C\|A\|\|B_y\|e^{v|t|-cd_\Lambda(a,y)}
        \label{eq:area_quasilocal_lr}
    \end{align}
    when $d_\Lambda(a,y)<\infty$, and the commutator is zero when
    $d_\Lambda(a,y)=\infty$. Moreover, for every integer $l\ge0$,
    \begin{align}
        \|\tau_t(A)-E_{N_l(a)}(\tau_t(A))\| & \le
        \|A\|\min\{2,Ce^{v|t|-cl}\}.
        \label{eq:area_quasilocal_lr_ce}
    \end{align}
    This is \cite[Lemma~4.1]{OpenAI2026AreaLaw}. Explicitly, for every
    $\mu\ge0$ and every integer $n\le d_\Lambda(a,y)$,
    $\|[\tau_t(A),B_y]\|\le2e^{\mu R}\|A\|\|B_y\|
        \exp(2v_R\mu_RJe^{2\mu R}|t|-\mu n)$.
\end{lemma}
\begin{proof}
    \leanok
    \uses{lem:area_domain_graph_counts,thm:qc_graph_lieb_robinson,
        thm:qc_graph_lieb_robinson_zero,thm:qc_graph_ball_localization}
    Apply Theorem~\ref{thm:qc_graph_lieb_robinson} and
    Lemma~\ref{thm:qc_graph_lieb_robinson_zero} on the induced graph of
    $\Lambda$, with an anchor family taking the value $a$ at $X$, support
    diameter $R$, support size $v_R$ and site budget $b_0=\mu_RJ$ from
    Lemma~\ref{lem:area_domain_graph_counts}. Apply
    Theorem~\ref{thm:qc_graph_ball_localization} with sphere growth
    $2(d+1)^2$. Both estimates have constants determined by $R$ and $J$.
    Take the larger of the two prefactors and of the two velocities, and the
    smaller of the two decay rates.
\end{proof}
